% 第2章 2.4.2：一致有界原理证明中的层级集示意（示意图，形状不代表具体空间）。
% F_k={x: sup_α‖A_α x‖≤k} 闭、对称、凸，且 F_k=kF_1；X=∪F_k。
% Baire 定理给出某个 F_{k0}（此处 k0=3）含球 B(x0,r)；把球平移到原点，
% 由 z=(x0+z)-x0 得 ‖z‖<r 时 ‖A_α z‖≤2k0，再缩放得 ‖A_α‖≤2k0/r。
\begin{tikzpicture}[x=10mm,y=10mm,font=\figurefont\footnotesize,text=FigureInk,
  note/.style={font=\FigureNoteFont\fontsize{7.5}{9}\selectfont,text=FigureMuted},
  vec/.style={line width=.9pt,-{Stealth[length=1.8mm]}}]
\begin{scope}[rotate=12]
  \fill[FigureBlueFill!60] (0,0) ellipse (3.3 and 1.35);
  \draw[FigureBlue,line width=.8pt] (0,0) ellipse (3.3 and 1.35);
  \draw[FigureBlue,line width=.8pt] (0,0) ellipse (2.2 and 0.9);
  \draw[FigureBlue,line width=.8pt] (0,0) ellipse (1.1 and 0.45);
  \coordinate (F1) at (0.85,-0.42); \coordinate (F2) at (1.75,-0.8); \coordinate (F3) at (2.6,-1.18);
  \coordinate (x0) at (-2.1,0.25);
\end{scope}
\node[text=FigureBlue,font=\figurefont\scriptsize] at (F1) {$F_1$};
\node[text=FigureBlue,font=\figurefont\scriptsize] at (F2) {$F_2$};
\node[text=FigureBlue,font=\figurefont\scriptsize] at (F3) {$F_3$};
\fill[FigureGreen!35] (x0) circle (0.38);
\draw[FigureGreen,line width=1pt] (x0) circle (0.38);
\fill[FigureInk] (x0) circle (.9pt) node[below=.3mm,font=\figurefont\scriptsize] {$x_0$};
\draw[FigureGreen,densely dashed,line width=.9pt] (0,0) circle (0.38);
\fill[FigureInk] (0,0) circle (.9pt) node[below=.4mm,font=\figurefont\scriptsize] {$0$};
\draw[vec,FigureStroke] ($(x0)+(0.15,0.42)$) to[bend left=25] node[above=1.2mm,note] {平移 $z=(x_0+z)-x_0$} (0.05,0.42);
\node[note,anchor=west] at (-3.4,-1.75) {$F_k=\{x:\sup_\alpha\|A_\alpha x\|\leq k\}=kF_1$，绿色球 $B(x_0,r)\subset F_3$};
\end{tikzpicture}
